\documentclass[11pt]{article}

\usepackage[margin=1in]{geometry}
\usepackage{amsmath,amssymb}
\usepackage{graphicx}
\usepackage[colorlinks=true,linkcolor=blue,citecolor=blue]{hyperref}
\usepackage[capitalize]{cleveref}
\usepackage{float}

% --- Supplementary numbering: figures/tables as S1, S2, ... (continuous, not per-section) ---
\renewcommand{\thefigure}{S\arabic{figure}}
\renewcommand{\thetable}{S\arabic{table}}
\renewcommand{\theequation}{S\arabic{equation}}

\title{Supplementary Material}
\author{}
\date{}

\begin{document}
\maketitle

\setcounter{page}{1}
\renewcommand{\thepage}{S\arabic{page}}

\section*{Supplementary Methods}

All analyses below use the same dual-reservoir pharmacokinetic/pharmacodynamic
model, treatment protocol, and baseline parameters as the main text
(vancomycin 1000~mg q8h for 4~days, followed by linezolid 600~mg q12h for
42~days, 21~days of pre-treatment equilibration and 14~days of post-treatment
follow-up; total simulated course 81~days). The limit of detection (LOD) is
10~CFU/mL in both compartments. Unless otherwise noted, growth rates are held
at their calibrated baseline values ($\rho_S = 0.61~\mathrm{h}^{-1}$,
$\rho_R = 0.55~\mathrm{h}^{-1}$, $\rho_S^r = 0.175~\mathrm{h}^{-1}$,
$\rho_R^r = 0.1765~\mathrm{h}^{-1}$), and only the single parameter under
study is varied. Escape thresholds are located by bisection
(\texttt{scipy.optimize.brentq}) on the difference between the final
simulated compartment count and the LOD, rather than read off a coarse grid,
so each reported boundary is precise to the stated tolerance.

% =============================================================================
\section{Deterministic escape threshold in linezolid potency ($EC_{50,L}$)}
\label{sec:s1}

\Cref{fig:ec50_threshold_sweep} shows the final $R_b$/$R_{res}$ burden across
a fine, deterministic sweep of $EC_{50,L} \in [0.3, 2.0]~\mathrm{mg/L}$
(171 points, step 0.01~mg/L), with every other parameter fixed at baseline.
Both compartments clear below the LOD for $EC_{50,L}$ below the bisected
threshold and persist above it, crossing together --- consistent with the
single shared escape boundary reported across all four sensitivity analyses
in the main text (\Cref{sec:s2} below reproduces this same threshold from the
Monte Carlo sweep). This confirms that the $EC_{50,L} \approx
0.932~\mathrm{mg/L}$ threshold quoted in the main text is a genuine, sharp
switch rather than an artifact of the coarser Monte Carlo sampling grid.

\begin{figure}[H]
    \centering
    \includegraphics[width=1.0\linewidth]{figures/ec50_lzd_threshold_sweep.png}
    \caption{Final $R_b$/$R_{res}$ burden versus $EC_{50,L}$ from the
    continuous deterministic sweep ($n = 171$ points, $EC_{50,L} \in
    [0.3, 2.0]~\mathrm{mg/L}$), with the bisected escape threshold marked for
    each compartment and the escape zone shaded.}
    \label{fig:ec50_threshold_sweep}
\end{figure}

% =============================================================================
\section{Monte Carlo escape switch in linezolid potency ($EC_{50,L}$)}
\label{sec:s2}

\Cref{fig:mc_ec50_switch} places the $n = 1000$ Monte Carlo patient
population sampled from $EC_{50,L} \sim \mathrm{LogNormal}$
(mean $1.0~\mathrm{mg/L}$, $\mathrm{CV} = 1.117$; the same sample used to
generate the $34.7\%$ escape fraction reported in the main text) against the
continuous deterministic switch curve from \Cref{sec:s1}. The top panel shows
final $R_b$/$R_{res}$ burden along a dense 121-point deterministic sweep of
the same range, with the bisected threshold marked; the bottom panel shows
where the sampled patient population actually falls relative to that switch,
with each patient colored by their own simulated outcome (escape vs.\
suppressed) rather than by which side of the threshold their sampled
$EC_{50,L}$ happens to fall on. The two panels agree exactly on where the
population divides, confirming that the Monte Carlo escape fraction reported
in the main text is driven entirely by this single deterministic boundary
rather than by any other source of between-run variability.

\begin{figure}[H]
    \centering
    \includegraphics[width=1.0\linewidth]{figures/mc_ec50_lzd_switch.png}
    \caption{Top: final $R_b$/$R_{res}$ burden versus $EC_{50,L}$ from the
    continuous deterministic sweep, with the bisected switch marked. Bottom:
    the $n = 1000$ Monte Carlo-sampled patient population, colored by each
    patient's own simulated outcome (escape vs.\ suppressed), split at the
    same threshold.}
    \label{fig:mc_ec50_switch}
\end{figure}

% =============================================================================
\section{Deterministic escape threshold in initial reservoir burden ($S_{res,0}$)}
\label{sec:s3}

\Cref{fig:sres0_threshold_sweep} sweeps the initial reservoir sensitive-strain
load $S_{res,0} \in [1, 2000]~\mathrm{CFU/mL}$ (101 points), holding
$R_{res,0} = 100~\mathrm{CFU/mL}$ and every other parameter at baseline, to
test how sensitive the escape outcome is to the size of the pre-treatment
reservoir population rather than to any pharmacodynamic parameter. Unlike the
$EC_{50,L}$ sweep in \Cref{sec:s1}, no clearance threshold exists anywhere in
the swept range: both $R_b$ and $R_{res}$ finish above the LOD at every one
of the 101 tested $S_{res,0}$ values, from 1 to 2000~CFU/mL. The model's own
default, $S_{res,0} = 100~\mathrm{CFU/mL}$, is marked directly on the panel
for reference, but is not distinguished from any other tested value in its
outcome. Resistant escape in this model is therefore not an artifact of the
specific initial reservoir burden assumed in the main text --- it persists
regardless of where $S_{res,0}$ is set within a two-decade range around that
default.

\begin{figure}[H]
    \centering
    \includegraphics[width=1.0\linewidth]{figures/sres0_threshold_sweep.png}
    \caption{Final $R_b$/$R_{res}$ burden versus initial reservoir sensitive
    load $S_{res,0}$, from a continuous deterministic sweep
    ($n = 101$ points, $S_{res,0} \in [1, 2000]~\mathrm{CFU/mL}$), with the
    model's own default ($S_{res,0} = 100$) marked. Both compartments remain
    above the LOD across the entire swept range; no clearance threshold
    exists in $S_{res,0}$ under baseline parameters.}
    \label{fig:sres0_threshold_sweep}
\end{figure}

% =============================================================================
\section{Scalar outcome sweeps for the compartment exchange rates}
\label{sec:s4}

\Cref{fig:supp_frb_outcomes,fig:supp_frb_outcomes_zoom} give the full-range
and zoomed scalar-outcome views for Sweep~1 ($f_{r \to b}$ varied, $f_{b \to
r}$ held at baseline). Peak $S_b$, peak $R_b$, and final $R_b$ all sit at or
near zero for $f_{r \to b} \lesssim 3\times10^{-5}~\mathrm{h}^{-1}$ --- blood
is seeded exclusively through reservoir-to-blood transfer, so vanishingly
small $f_{r \to b}$ leaves bacteremia unestablished for the entire simulated
course --- then switch sharply to their full baseline levels within roughly
half a log-decade of the model's own baseline rate
($5\times10^{-5}~\mathrm{h}^{-1}$). Peak $S_{res}$ is flat across the full
six-decade range. At the upper end, final $R_b$ undergoes a second,
independent transition, dropping back to zero for
$f_{r \to b} \gtrsim 5\times10^{-3}~\mathrm{h}^{-1}$ even as peak $R_b$
continues to rise --- $R_b$ still establishes and peaks at these rates, but no
longer relapses by the end of follow-up.

\begin{figure}[H]
    \centering
    \includegraphics[width=1.0\linewidth]{figures/mc_frb_sweep_outcomes.png}
    \caption{Scalar outcomes (peak $S_b$, peak $R_b$, peak $S_{res}$, final
    $R_b$) versus $f_{r \to b}$ across the full log-spaced grid,
    $f_{b \to r}$ held at baseline.}
    \label{fig:supp_frb_outcomes}
\end{figure}

\begin{figure}[H]
    \centering
    \includegraphics[width=1.0\linewidth]{figures/mc_frb_sweep_outcomes_zoom.png}
    \caption{Scalar outcomes versus $f_{r \to b}$, restricted to the
    $[1\times10^{-3}, 1\times10^{-2}]~\mathrm{h}^{-1}$ transition band.}
    \label{fig:supp_frb_outcomes_zoom}
\end{figure}

\Cref{fig:supp_fbr_outcomes,fig:supp_fbr_outcomes_zoom} give the same two
views for Sweep~2 ($f_{b \to r}$ varied, $f_{r \to b}$ held at baseline).
Peak $S_b$, peak $R_b$, and final $R_b$ are flat to within 1\% across roughly
the lower four and a half of the six log-decades sampled
($10^{-8}$--$10^{-3.5}~\mathrm{h}^{-1}$), consistent with $f_{b \to r}$
acting only as an efflux from blood rather than a seeding route for either
compartment. Peak $S_{res}$ is not flat over the upper range, however: it
rises from its baseline value ($\sim$4560~CFU/mL) to a maximum of
$\sim$6270~CFU/mL ($+37\%$) near $f_{b \to r} \approx
3\times10^{-3}~\mathrm{h}^{-1}$, before collapsing back toward baseline as
peak $S_b$ simultaneously collapses toward zero, both by
$f_{b \to r} = 1\times10^{-2}~\mathrm{h}^{-1}$.

\begin{figure}[H]
    \centering
    \includegraphics[width=1.0\linewidth]{figures/mc_fbr_sweep_outcomes.png}
    \caption{Scalar outcomes (peak $S_b$, peak $R_b$, peak $S_{res}$, final
    $R_b$) versus $f_{b \to r}$ across the full log-spaced grid,
    $f_{r \to b}$ held at baseline.}
    \label{fig:supp_fbr_outcomes}
\end{figure}

\begin{figure}[H]
    \centering
    \includegraphics[width=1.0\linewidth]{figures/mc_fbr_sweep_outcomes_zoom.png}
    \caption{Scalar outcomes versus $f_{b \to r}$, restricted to the
    $[1\times10^{-3}, 1\times10^{-2}]~\mathrm{h}^{-1}$ transition band.}
    \label{fig:supp_fbr_outcomes_zoom}
\end{figure}

% =============================================================================
\section{Sampling grid for the compartment exchange-rate sweeps}
\label{sec:s5}

The two exchange-rate sweeps in \Cref{sec:s4}
($f_{r \to b}$, reservoir $\to$ blood; $f_{b \to r}$, blood $\to$ reservoir)
each step their target rate across a shared log-spaced grid spanning six
decades ($10^{-8}$--$10^{-2}~\mathrm{h}^{-1}$), while holding the other
exchange rate fixed at its own baseline value
($f_{r \to b} = 5\times10^{-5}~\mathrm{h}^{-1}$,
$f_{b \to r} = 1\times10^{-5}~\mathrm{h}^{-1}$). Because both transition bands
identified in the main text fall within $[10^{-3}, 10^{-2}]~\mathrm{h}^{-1}$,
the grid is deliberately non-uniform: a coarse 40-point log-spaced pass covers
the full six-decade range for context, supplemented by a dense 160-point pass
restricted to that transition band, so that the sharp switches reported in
\Cref{fig:supp_frb_outcomes_zoom,fig:supp_fbr_outcomes_zoom} are resolved
rather than aliased. \Cref{fig:transfer_rate_grid} shows the resulting grid
coverage for both sweeps directly, with each sweep's own baseline rate
marked.

\begin{figure}[H]
    \centering
    \includegraphics[width=1.0\linewidth]{figures/mc_transfer_rate_sweep_grid.png}
    \caption{Log-spaced grid coverage used for the two exchange-rate sweeps
    reported in the main text ($n = 193$ points per sweep). Top: grid for
    Sweep~1 ($f_{r \to b}$ varied, $f_{b \to r}$ held at baseline). Bottom:
    grid for Sweep~2 ($f_{b \to r}$ varied, $f_{r \to b}$ held at baseline).
    Each sweep concentrates 160 of its 193 points inside the
    $[10^{-3}, 10^{-2}]~\mathrm{h}^{-1}$ transition band identified in the
    main text, with the remaining 33 points spread log-uniformly across the
    full six-decade range for context. The band's boundaries are not
    arbitrary: $1/(10^{-3}~\mathrm{h}^{-1}) \approx 42$~days and
    $1/(10^{-2}~\mathrm{h}^{-1}) \approx 4$~days match the linezolid
    (42-day) and vancomycin (4-day) course lengths almost exactly, so this
    band brackets the regime where the exchange timescale is comparable to a
    treatment-phase duration -- below $10^{-3}~\mathrm{h}^{-1}$ exchange is
    too slow to complete even once during the linezolid course, and above
    $10^{-2}~\mathrm{h}^{-1}$ the compartments equilibrate within the
    (shorter) vancomycin window, so outcomes saturate on either side.}
    \label{fig:transfer_rate_grid}
\end{figure}

% =============================================================================
\section{Deterministic escape threshold in host immune clearance ($k_{immune}$)}
\label{sec:s9}

\Cref{fig:immune_threshold_sweep} shows the final $R_b$/$R_{res}$ burden
across a fine, deterministic sweep of $k_{immune} \in
[0.02, 0.22]~\mathrm{h}^{-1}$ (101 points), with every other parameter fixed
at baseline. As with $E_{max,L}$, escape occurs \emph{below} the threshold
rather than above it: $k_{immune}$ is a clearance rate, so weaker immune
clearance favors persistence. Both compartments clear above the bisected
threshold and persist below it, crossing together --- consistent with the
single shared escape boundary reported across all four sensitivity analyses
in the main text. The model's own default, $k_{immune} =
0.12~\mathrm{h}^{-1}$, is marked directly on the panel and falls well within
the escape zone (\Cref{sec:s6} below reproduces this same threshold from the
Monte Carlo sweep).

\begin{figure}[H]
    \centering
    \includegraphics[width=1.0\linewidth]{figures/immune_response_threshold_sweep.png}
    \caption{Final $R_b$/$R_{res}$ burden versus $k_{immune}$ from the
    continuous deterministic sweep ($n = 101$ points, $k_{immune} \in
    [0.02, 0.22]~\mathrm{h}^{-1}$), with the bisected escape threshold and
    the model's own default ($k_{immune} = 0.12~\mathrm{h}^{-1}$) marked.}
    \label{fig:immune_threshold_sweep}
\end{figure}

% =============================================================================
\section{Monte Carlo escape switch in host immune clearance ($k_{immune}$)}
\label{sec:s6}

\Cref{fig:supp_immune_switch} places the $n = 1000$ Monte Carlo patient
population sampled from $k_{immune} \sim \mathrm{LogNormal}$
(mean $0.17~\mathrm{h}^{-1}$, $\mathrm{CV} = 0.25$; the same sample used to
generate the $35.7\%$ relapse rate reported in the main text) against the
continuous deterministic switch curve from \Cref{sec:s9}. As with
$EC_{50,L}$ in \Cref{sec:s2}, the top panel
shows the continuous deterministic switch and the bottom panel shows where
the sampled patient population falls relative to it, each patient colored by
its own simulated outcome rather than by which side of the threshold its
sampled $k_{immune}$ happens to fall on. Because the bisected threshold sits
below the sampled mean, most simulated patients draw a $k_{immune}$ value
sufficient to clear the resistant subpopulations --- but with
$\mathrm{CV} = 0.25$, just over a third draw a value too low to do so,
confirming that resistant-strain persistence remains meaningfully sensitive
to inter-patient variability in host immune clearance even though the
majority of the population resolves the infection.

\begin{figure}[H]
    \centering
    \includegraphics[width=1.0\linewidth]{figures/mc_immune_response_switch.png}
    \caption{Top: final $R_b$/$R_{res}$ burden versus $k_{immune}$ from the
    continuous deterministic sweep, with the bisected switch marked. Bottom:
    the $n = 1000$ Monte Carlo-sampled patient population, colored by each
    patient's own simulated outcome (escape vs.\ suppressed), split at the
    same threshold.}
    \label{fig:supp_immune_switch}
\end{figure}

% =============================================================================
\section{Deterministic escape threshold in linezolid killing capacity ($E_{max,L}$)}
\label{sec:s7}

\Cref{fig:emax_threshold_sweep} shows the final $R_b$/$R_{res}$ burden across
a fine, deterministic sweep of $E_{max,L} \in [0.0, 1.0]~\mathrm{h}^{-1}$
(101 points, step 0.01), with every other parameter fixed at baseline.
Unlike $EC_{50,L}$, escape here occurs \emph{below} the threshold rather than
above it: $E_{max,L}$ is linezolid's ceiling effect, so a weaker ceiling
fails to hold resistant growth in check. Both compartments clear above the
bisected threshold and persist below it, crossing together --- consistent
with the single shared escape boundary reported across all four sensitivity
analyses in the main text (\Cref{sec:s8} below reproduces this same
threshold from the Monte Carlo sweep).

\begin{figure}[H]
    \centering
    \includegraphics[width=1.0\linewidth]{figures/emax_lzd_threshold_sweep.png}
    \caption{Final $R_b$/$R_{res}$ burden versus $E_{max,L}$ from the
    continuous deterministic sweep ($n = 101$ points, $E_{max,L} \in
    [0.0, 1.0]~\mathrm{h}^{-1}$), with the bisected escape threshold marked
    for each compartment and the escape zone shaded.}
    \label{fig:emax_threshold_sweep}
\end{figure}

% =============================================================================
\section{Monte Carlo escape switch in linezolid killing capacity ($E_{max,L}$)}
\label{sec:s8}

\Cref{fig:supp_emax_switch} places the $n = 1000$ Monte Carlo patient
population sampled from $E_{max,L} \sim \mathrm{LogNormal}$
(mean $0.8~\mathrm{h}^{-1}$, $\mathrm{CV} = 0.30$; the same sample used to
generate the $56.5\%$ escape fraction reported in the main text) against the
continuous deterministic switch curve from \Cref{sec:s7}. As with
$EC_{50,L}$ in \Cref{sec:s2}, the top panel shows the continuous switch and
the bottom panel shows where the sampled patient population falls relative
to it, each patient colored by its own simulated outcome rather than by
which side of the threshold its sampled $E_{max,L}$ happens to fall on.
Because the bisected threshold ($0.8161~\mathrm{h}^{-1}$) sits just above
the sampled mean, just over half of simulated patients drew an $E_{max,L}$
value too low to clear the infection, and the population divides exactly at
the same threshold located by the deterministic sweep.

\begin{figure}[H]
    \centering
    \includegraphics[width=1.0\linewidth]{figures/mc_emax_lzd_switch.png}
    \caption{Top: final $R_b$/$R_{res}$ burden versus $E_{max,L}$ from the
    continuous deterministic sweep, with the bisected switch marked. Bottom:
    the $n = 1000$ Monte Carlo-sampled patient population, colored by each
    patient's own simulated outcome (escape vs.\ suppressed), split at the
    same threshold.}
    \label{fig:supp_emax_switch}
\end{figure}

\end{document}
